\section{Marco Teórico}

\subsection{Computación de Alto Rendimiento (HPC)}
La Computación de Alto Rendimiento (HPC, por sus siglas en inglés) consiste en la agregación de potencia de cálculo a través de arquitecturas paralelas para resolver problemas científicos y tecnológicos de extrema complejidad. Esta disciplina permite ejecutar algoritmos y modelados matemáticos que resultarían inabarcables para un sistema tradicional debido a restricciones de tiempo y recursos de hardware \parencite{Hager2010HPC}. Al dividir cargas masivas de trabajo en instrucciones más pequeñas que se procesan simultáneamente, HPC maximiza la velocidad de procesamiento de datos y la obtención de resultados en la ingeniería \parencite{Patterson2017Computer}.

\subsection{Fundamentos Matemáticos de la Multiplicación Matricial}
La multiplicación de matrices es una operación algebraica bidimensional fundamental que procesa dos matrices de origen para producir una tercera. Dadas dos matrices \(A\) de dimensión \(m \times p\) y \(B\) de dimensión \(p \times n\), su producto da como resultado una matriz \(C\) de dimensiones \(m \times n\). Matemáticamente, el valor de cada celda individual \(C_{ij}\) en la matriz resultante se calcula mediante el producto punto o escalar entre la \(i\)-ésima fila de la matriz \(A\) y la \(j\)-ésima columna de la matriz \(B\). Esta operación se expresa de forma analítica mediante la siguiente sumatoria:

\[ C_{ij} = \sum_{k=1}^{p} A_{ik} B_{kj} \]

Para que esta operación esté matemáticamente definida y sea computacionalmente viable, existe una condición estricta de dimensionalidad: el número de columnas de la primera matriz (\(p\)) debe coincidir exactamente con el número de filas de la segunda matriz \parencite{Strang2016Linear}. 

En el contexto de la ingeniería y la Computación de Alto Rendimiento (HPC), la multiplicación matricial trasciende la teoría pura para consolidarse como la rutina principal (\textit{kernel}) subyacente en la gran mayoría del software científico. Esta operación constituye el pilar del Nivel 3 de la especificación BLAS (\textit{Basic Linear Algebra Subprograms}), el estándar de facto para operaciones de álgebra lineal. La eficiencia con la que un procesador ejecuta este algoritmo dicta directamente el rendimiento de aplicaciones críticas modernas, tales como el modelado de dinámica de fluidos, la física de partículas, el procesamiento masivo de imágenes y, muy especialmente, el entrenamiento de arquitecturas profundas en inteligencia artificial mediante redes neuronales \parencite{Dongarra1990BLAS}.


\subsection{Complejidad Computacional de la Multiplicación de Matrices}
La multiplicación de matrices es un problema canónico en el álgebra lineal computacional, utilizado frecuentemente para evaluar el desempeño y la robustez de las arquitecturas de hardware. Dadas dos matrices cuadradas \(A\) y \(B\) de dimensiones \(n \times n\), el algoritmo iterativo secuencial clásico exige iterar mediante tres ciclos anidados, realizando matemáticamente \(n^3\) multiplicaciones y \(n^2(n-1)\) sumas \parencite{Cormen2009Algorithms}. 

Esto establece una complejidad temporal asintótica estricta de \(\mathcal{O}(n^3)\). Este crecimiento cúbico implica que, si el tamaño de la matriz se duplica, el esfuerzo computacional se multiplica por un factor de ocho. Debido a esta pronunciada curva de costo algorítmico, la distribución de la carga mediante técnicas de programación paralela (hilos o procesos) se vuelve un requisito indispensable cuando se operan matrices de dimensiones elevadas \parencite{Grama2003Introduction}.

\subsection{Programación Secuencial y Concurrente}
La programación secuencial es el paradigma clásico de desarrollo donde las instrucciones de un algoritmo se ejecutan de forma estrictamente lineal, una tras otra, sobre un único núcleo de procesamiento. En este modelo tradicional, el límite de rendimiento está directamente dictado por la frecuencia de reloj del hardware. Como respuesta a estas limitaciones físicas, la programación concurrente estructura el software en múltiples hilos o tareas independientes que hacen progreso de forma superpuesta en el tiempo. Cuando estas tareas operan de manera simultánea en una arquitectura de hardware multinúcleo, la concurrencia se traduce en paralelismo físico, lo cual permite dividir la carga de trabajo y reducir drásticamente el tiempo de ejecución en problemas de cómputo intensivo \parencite{BenAri2006Concurrent}.

\subsection{Sistemas de Memoria Compartida}
En el ámbito de la computación paralela, el modelo de memoria compartida establece que múltiples hilos de ejecución (como los creados por la librería Pthreads) operan dentro de un mismo espacio de direccionamiento lógico. A nivel de arquitectura, esto significa que todos los núcleos del procesador pueden acceder de forma directa y simultánea a las mismas estructuras de datos residentes en la memoria principal. Para problemas matriciales, este modelo resulta excepcionalmente eficiente frente a los sistemas de memoria distribuida, ya que evita la necesidad de copiar y transmitir los datos de las matrices de origen repetidamente entre procesos. Aunque este paradigma exige control riguroso para evitar condiciones de carrera, en operaciones intrínsecamente independientes como el cálculo individual de las celdas de una matriz resultado, su implementación maximiza la eficiencia computacional sin sobrecarga de sincronización \parencite{Silberschatz2018Operating}.

\subsection{Jerarquía de Caché, Líneas de Caché y Localidad de Memoria}
El rendimiento de una aplicación de alto rendimiento (HPC) no depende exclusivamente de la capacidad matemática del procesador, sino también del acceso al subsistema de memoria. Para mitigar la alta latencia térmica y física de la memoria RAM, los procesadores modernos integran una jerarquía de memoria ultrarrápida cercana a los núcleos, conocida como memoria Caché, habitualmente dividida en niveles (L1, L2 y L3). La unidad mínima de transferencia de información entre la RAM principal y la memoria caché no es un byte aislado, sino un bloque continuo de datos denominado línea de caché (\textit{cache line}), que en las arquitecturas modernas suele constar de 64 bytes \parencite{Drepper2007Memory}.

Esta transferencia por bloques está fundamentada en el principio de localidad espacial de memoria, un concepto empírico que dicta que si un programa accede a una dirección de memoria, es inminentemente probable que acceda a las direcciones adyacentes a continuación. En lenguajes de programación como C, las matrices bidimensionales se almacenan en la memoria RAM de manera contigua organizadas por filas (\textit{row-major order}). Durante la multiplicación matricial clásica iterativa, la segunda matriz se recorre inevitablemente por columnas; este patrón de acceso ortogonal rompe por completo la localidad espacial. Al intentar leer datos saltando grandes bloques de memoria, los valores necesarios no se encuentran en la línea de caché cargada, provocando fallos de caché (\textit{cache misses}) que obligan al procesador a detenerse (\textit{stall}) mientras busca los datos nuevamente en la lenta memoria RAM, lo cual penaliza severamente el desempeño del algoritmo \parencite{Patterson2017Computer}.

\subsection{Optimización Algorítmica y Aprovechamiento de la Caché}
Aunque la paralelización distribuye la carga de trabajo, el cuello de botella de la memoria solo puede resolverse reestructurando el algoritmo subyacente. La estrategia de optimización más directa para la multiplicación de matrices en C consiste en el intercambio de bucles (\textit{loop interchange}). Altera el orden tradicional de los ciclos anidados de \(i, j, k\) a \(i, k, j\). Matemáticamente, el resultado es idéntico, pero a nivel de hardware, el bucle más interno ahora recorre las matrices \(B\) y \(C\) estrictamente por filas. Esto garantiza que el procesador consuma por completo los 64 bytes de cada línea de caché (\textit{cache line}) cargada antes de solicitar la siguiente, aprovechando al máximo la localidad espacial y reduciendo las penalizaciones por fallos de caché en órdenes de magnitud \parencite{Lam1991Cache}.

Para matrices masivas que exceden la capacidad total de la memoria caché (L1 y L2), la literatura de HPC emplea una técnica avanzada denominada multiplicación por bloques o \textit{Loop Tiling}. Este método particiona las matrices originales en submatrices (o bloques) de un tamaño cuidadosamente calculado para que residan por completo dentro de la caché rápida del procesador. En lugar de multiplicar filas y columnas completas, el algoritmo opera sobre estos bloques aislados, reutilizando los datos cargados repetidas veces antes de que el controlador de memoria los descarte (\textit{eviction}). Esta técnica transforma un problema limitado por el ancho de banda de la memoria (\textit{memory-bound}) en uno puramente computacional (\textit{compute-bound}) \parencite{Grama2003Introduction}.

\subsection{Métricas de Rendimiento en Computación Secuencial}
El análisis de un algoritmo secuencial requiere establecer una línea base absoluta sobre la cual evaluar cualquier mejora algorítmica o de hardware. La métrica fundamental empírica es el Tiempo de Ejecución (\textit{Wall-clock time}), que representa el tiempo físico real transcurrido desde el inicio hasta la finalización del programa. Sin embargo, dado que este valor fluctúa dependiendo de la interferencia del sistema operativo y los accesos a memoria, la arquitectura de computadores recurre a métricas de rendimiento intrínseco (\textit{throughput}) \parencite{Patterson2017Computer}.

Para cargas de trabajo de alto rendimiento algebraico como la multiplicación de matrices, la métrica estandarizada predominante es el volumen de Operaciones de Punto Flotante por Segundo (FLOPS, por sus siglas en inglés). Esta medida cuantifica el rendimiento matemático bruto del procesador y se define mediante la siguiente relación:

\begin{equation}
    \text{FLOPS} = \frac{\text{Operaciones Aritméticas Totales}}{T_{\text{ejecución}}}
\end{equation}

Para el algoritmo clásico iterativo que multiplica dos matrices cuadradas de dimensión \(n \times n\), el número total de operaciones de punto flotante se establece analíticamente en \(2n^3\) (contabilizando una instrucción de multiplicación y una de suma por cada iteración del bucle más interno). Al contrastar los FLOPS experimentales obtenidos en la ejecución secuencial contra el rendimiento teórico máximo (\textit{Peak FLOPS}) que el fabricante del procesador garantiza para un solo núcleo, es posible diagnosticar ineficiencias de memoria o cuellos de botella antes de recurrir a la paralelización \parencite{Hager2010HPC}. 

Adicionalmente, el rendimiento de la ejecución en un único núcleo se evalúa a nivel microarquitectónico mediante el IPC (Instrucciones por Ciclo de reloj). Esta métrica indica la eficiencia con la que el código compilado logra mantener ocupada la segmentación (\textit{pipeline}) del procesador, minimizando las detenciones o burbujas producidas por dependencias de datos en la ejecución lineal \parencite{Yi2020CPI}.


\subsection{Computación paralela}

La computación paralela es un paradigma de procesamiento en el que un problema de gran tamaño se divide en subproblemas más pequeños que pueden resolverse de forma simultánea mediante múltiples procesadores o núcleos de ejecución. Su objetivo principal es reducir el tiempo total de cómputo y aumentar la eficiencia de uso de los recursos disponibles, especialmente en problemas con alta carga aritmética o grandes volúmenes de datos \parencite{ibmParallelComputing2024}.

Desde una perspectiva conceptual, la computación paralela contrasta con la computación secuencial, en la cual las instrucciones se ejecutan una tras otra sobre un único procesador. En entornos paralelos, el rendimiento no depende únicamente de la capacidad de cálculo del procesador, sino también de la forma en que se particiona el trabajo, se distribuyen los datos y se coordinan los procesos durante la ejecución \parencite{grama2003introduction}.

En la práctica, la computación paralela puede implementarse sobre diferentes arquitecturas, entre ellas memoria compartida, memoria distribuida y esquemas híbridos. En sistemas de memoria distribuida, cada proceso posee su propia memoria local y la coordinación entre nodos se realiza mediante intercambio explícito de mensajes. Este modelo resulta especialmente útil en clústeres y supercomputadores, donde la escalabilidad depende de minimizar la sobrecarga de comunicación y maximizar la proporción entre cómputo útil y costo de coordinación \parencite{hager2010introHPC}.

La relevancia de la computación paralela en computación de alto rendimiento proviene de su capacidad para resolver problemas que serían impracticables en forma secuencial. Aplicaciones como simulación numérica, álgebra lineal densa, análisis de datos y aprendizaje automático se benefician de este paradigma debido a que presentan un alto grado de paralelismo natural y un costo computacional considerable \parencite{patterson2017computer}.

\subsection{MPI}

MPI (\textit{Message Passing Interface}) es un estándar para programación paralela mediante paso de mensajes, diseñado principalmente para sistemas de memoria distribuida. Su propósito es ofrecer una interfaz portable, eficiente y ampliamente soportada para desarrollar aplicaciones paralelas en lenguajes como C, C++ y Fortran \parencite{snir1998mpi2,mpiForum2021}.

Una de las características fundamentales de MPI es que el programador controla explícitamente la descomposición del problema, la distribución de datos y la comunicación entre procesos. Esto permite construir aplicaciones con un modelo SPMD (\textit{Single Program, Multiple Data}), en el cual múltiples procesos ejecutan el mismo programa pero trabajan sobre datos distintos \parencite{gropp2014usingmpi}.

En MPI, los procesos se organizan dentro de comunicadores, que definen conjuntos de procesos autorizados para intercambiar mensajes. Cada proceso dentro de un comunicador posee un rango (\textit{rank}) que lo identifica de forma única y permite establecer el origen o destino de la comunicación. Esta abstracción facilita tanto las comunicaciones punto a punto como las operaciones colectivas.

Entre las operaciones más utilizadas de MPI se encuentran el envío y recepción de mensajes entre procesos, así como las comunicaciones colectivas como difusión (\textit{broadcast}), dispersión (\textit{scatter}), reunión (\textit{gather}) y reducción (\textit{reduce}). Estas primitivas permiten implementar de manera estructurada algoritmos paralelos donde los datos deben repartirse, procesarse localmente y recombinarse al final de la ejecución \parencite{mpiForum2021}.

En el contexto de la multiplicación de matrices, MPI es particularmente adecuado porque el problema puede descomponerse por filas o bloques, de manera que cada proceso calcula una parte independiente de la matriz resultado. No obstante, el rendimiento final depende de la relación entre el tamaño del problema, el costo de comunicación y la eficiencia del núcleo de cómputo local. Por esta razón, las implementaciones MPI suelen requerir además optimización algorítmica y cuidadosa atención al acceso a memoria \parencite{snir1998mpi2,pacheco1997parallel}.

\subsection{Network File System (NFS)}

El Network File System (NFS) es un protocolo de sistema de archivos distribuido que permite a varios nodos acceder a archivos almacenados de forma remota como si estuvieran disponibles localmente en el sistema operativo. Su principal ventaja es que proporciona un mecanismo transparente de compartición de datos, de manera que los procesos pueden leer y escribir archivos ubicados en un servidor central sin requerir una lógica explícita de transferencia de archivos en la aplicación \parencite{sandberg1985nfs,stoner1995design}.

En entornos de computación paralela y distribuida, NFS resulta útil para compartir binarios, datos de entrada y resultados entre múltiples nodos de un clúster. Esta característica simplifica el despliegue de aplicaciones MPI, ya que todos los procesos pueden acceder a la misma estructura de directorios y ejecutar el mismo binario sin necesidad de copiar manualmente los archivos en cada máquina \parencite{tanenbaum2015distributed,storm1988nfs}.

Una propiedad importante de NFS es que el acceso a archivos remotos introduce costos adicionales de red y de sincronización frente a un sistema de archivos local. Aunque el cliente percibe la operación como una lectura o escritura convencional, la solicitud debe transmitirse al servidor y esperar respuesta, lo que puede afectar el tiempo total de ejecución si el volumen de accesos es alto o si el sistema se utiliza simultáneamente por varios procesos \parencite{tanenbaum2015distributed,stone1995filesystem}.

En un clúster MPI, este comportamiento tiene implicaciones directas sobre el rendimiento. Si los archivos de entrada se leen desde NFS al inicio de la ejecución, la sobrecarga puede ser aceptable; sin embargo, si la aplicación depende de accesos frecuentes o si varios procesos realizan lecturas concurrentes sobre archivos grandes, la latencia de red y el ancho de banda del sistema de almacenamiento pueden convertirse en un cuello de botella. Por esta razón, el uso de NFS debe analizarse no solo como una comodidad operativa, sino también como un factor que influye en las mediciones experimentales \parencite{stone1995filesystem,terry1995session}.

En este proyecto, NFS cumple una función de infraestructura para distribuir los binarios, leer las matrices de entrada y almacenar resultados en una ubicación común. Esto facilita la ejecución homogénea en todos los nodos y evita inconsistencias entre entornos locales, pero al mismo tiempo introduce una capa adicional de dependencia sobre el subsistema de almacenamiento compartido, que debe considerarse al interpretar los tiempos medidos \parencite{sandberg1985nfs,tanenbaum2015distributed}.